\section{问题三的模型建立与求解}

\subsection{通信服务状态模型}

问题三在问题二运输调度基础上叠加连续通信约束。运输无人机在爬升、巡航、下降及投送全程须保持与固定网关 $G01$ 的直连；直连不可用时，可由一架中继无人机提供“运输—中继—$G01$”两段链路。依据附录 3 与自由空间衰减标准\cite{ref_itu,ref_zeng}，端点 $a,b$ 在三维距离 $d_{ab}$（km）下的自由空间损耗为
\begin{equation}
L_{fs}=32.44+20\log_{10} d_{ab}+20\log_{10} f,\qquad f=2400\,\text{MHz},
\end{equation}
考虑地形遮挡后总损耗 $L=L_{fs}+\mathbb{1}[\text{遮挡}]\cdot L_{obs}$，其中遮挡由两端点视线连线与 $30\,$m DEM 的相交判定，$L_{obs}=10\,$dB。方向 $a\to b$ 的最大允许损耗
\begin{equation}
L^{\max}_{a\to b}=P_t^{a}+G^{a}+G^{b}-L_{sys}-\big(P_{sens}^{b}+M\big),
\end{equation}
双向链路门限取两方向较小值 $L^{\text{th}}_{ab}=\min\{L^{\max}_{a\to b},L^{\max}_{b\to a}\}$。链路可用当且仅当 $L\le L^{\text{th}}_{ab}$。运输无人机某时刻通信状态判定为：直连可用记直连；否则若接入链路（运输—中继）与回传链路（中继—$G01$）\emph{同时}可用记中继；其余记中断。

\subsection{中继悬停网格与联合调度模型}

将每条运输航段按时间/空间采样为若干通信相位，逐相位判定直连是否可用。统计得需中继保障的相位 $586/720$。对存在缺口的航段，将中继悬停位置离散化为候选网格：沿缺口腿水平连线取 $5$ 个分数点，每点配 $6$ 档悬停离地高度（$50\!\sim\!300\,$m），共 $3420$ 个候选悬停点；零候选（硬缺口）相位数为 $0$，说明候选网格足以覆盖全部缺口。

中继无人机由 $O01$ 出发，飞抵悬停点建链服务后返回，其飞行时间/能耗按附录 2 三段式规则以计划起飞总质量计算，服务能耗由悬停功率与通信附加功率确定，并满足返航电量要求。联合调度在问题二决策变量之外新增：中继无人机悬停位置、飞行海拔、服务时段、能源组件分配。目标扩展为运输与中继两类无人机的加权迟到量、联合完工时间、总能耗（运输＋中继）与总架次：
\begin{equation}
\min\ \big(f_1^{\text{late}},\ f_2^{\text{makespan}},\ f_3^{E_{\text{tr}}+E_{\text{re}}},\ f_4^{N_{\text{tr}}+N_{\text{re}}}\big).
\end{equation}

\subsection{求解方法与结果}

\textbf{高复杂度主解（全局分组编码 NSGA-II）.} 在问题二的全局分组编码基础上，路线几何先确定每条运输航段的通信缺口，再由候选悬停点池和联合离散事件仿真完成中继覆盖与资源调度；中继位置不是额外的自由搜索基因。正式种群 $80$、代数 $80$，总耗时 $672.6\,$s。代表解：$N=90$（运输 $55$＋中继 $35$），联合完工 $f_2=54352.2\,$s，总能耗 $f_3=135.639\,$kWh（运输 $110.856$＋中继 $24.782$），加权迟到量 $f_1=7.49\times10^{6}$，硬约束违反 $12$；$55$ 条运输路线中直连 $20$、中继覆盖 $35$、通信缺口 $0$，$80/80$ 个货箱送达。

\textbf{基础解（贪心运输路线＋中继覆盖分配）.} 承接问题二贪心方案（$16$ 个运输架次），逐相位判定直连，并从候选网格中贪心分配中继悬停点。结果为总架次 $29$（运输 $16$＋中继 $13$）、$f_2=19379.4\,$s、$f_3=67.387\,$kWh，硬约束违反数 $18$；其中 $16$ 条运输路线直连 $3$、中继覆盖 $13$、通信缺口 $0$。该方案的作用是提供透明、低复杂度的对照，不与主解直接竞争全局最优性。

通信保障链路预算、悬停网格、联合调度时序与两方法对照如图~\ref{fig:q3} 所示。两条方案的数据链不同：基础解直接继承问题二贪心路线，其运输能耗仍为 $58.179\,$kWh，新增中继能耗为 $9.208\,$kWh；高复杂度解则重新搜索运输分组，运输能耗为 $110.856\,$kWh，高于问题二全局 NSGA-II 的 $82.777\,$kWh，另需中继能耗 $24.782\,$kWh。因此表中差异反映端到端方案及其路线重构，不应将总能耗增量解释为单纯通信保障代价。两方案均实现\textbf{通信缺口为零}，验证了“缺口腿分数点×悬停高度档”候选网格设计在本次计算场景下的覆盖效果。

\begin{table}[htbp]
\centering
\caption{问题三高复杂度主解与基础解对比}
\label{tab:q3_methods}
\renewcommand{\arraystretch}{1.2}
\small
\begin{tabular}{c|c|c|c|c|c|c}
\toprule
方案 & 运输+中继架次 & 总能耗/kWh & 联合完工/s & 直连/中继 & 通信缺口 & 违反数 \\
\midrule
高复杂度主解 & $55+35=90$ & 135.639 & 54352.2 & 20/35 & 0 & 12 \\
基础解 & $16+13=29$ & 67.387 & 19379.4 & 3/13 & 0 & 18 \\
\bottomrule
\end{tabular}
\end{table}

\begin{figure}[htbp]
  \centering
  \subfloat[双向链路预算与直连可用性]{\includegraphics[width=0.48\textwidth]{q3_1_01.png}}
  \hfill
  \subfloat[中继悬停候选网格与覆盖]{\includegraphics[width=0.48\textwidth]{q3_1_02.png}}
  \\
  \subfloat[运输—中继联合调度时序]{\includegraphics[width=0.48\textwidth]{q3_1_03.png}}
  \hfill
  \subfloat[两方法能耗与架次对照]{\includegraphics[width=0.48\textwidth]{q3_2_01.png}}
  \caption{问题三通信约束下联合调度结果}
  \label{fig:q3}
\end{figure}

\subsection{通信保障与中继能耗的多维分析}

对 NSGA-II 联合优化解的通信结构与中继能耗进一步分析如图~\ref{fig:adv_q3}。图~\ref{fig:adv_q3}(a) 的环形图显示 $55$ 条运输路线中直连 $20$、需中继保障 $35$，中继占比约 $64\%$，说明山区地形遮挡下多数航段无法维持直连、通信中继是刚需。图~\ref{fig:adv_q3}(b) 按通信方式区分路线能耗—时长关系，可见直连路线（绿色）多为近端短航时低能耗，需中继的路线（橙色）则分布在更长航时、更高能耗区间，二者随航时同步上升。图~\ref{fig:adv_q3}(c) 的中继悬停高度直方图呈现双峰：低空（约 $100\,$m）用于近距补盲、高空（约 $300\,$m）用于跨越山脊，反映候选网格“高度分档”设计被充分利用。图~\ref{fig:adv_q3}(d) 表明中继架次能耗随在位服务时长近似线性增长，且高 AGL（深色）点整体偏上，即“飞得高、守得久”是中继能耗的主要来源。

\begin{figure}[htbp]
  \centering
  \includegraphics[width=0.96\textwidth]{adv_q3.png}
  \caption{通信保障与中继能耗分析}
  \label{fig:adv_q3}
\end{figure}
